\documentclass[10pt,a4paper]{amsart}
\usepackage{amsmath,amssymb,amsthm,graphicx}
% Retained minimal preamble for the Stacks Project chapter flat.tex (More on Flatness).
% Verbatim excerpt of the chapter preamble: Xy-pic setup (native preamble lines 14-18) and the native
% theorem declarations with the mathematical macros (native preamble lines 154-217). Copied mechanically;
% nothing added, removed or redefined. The wrapper class and the amsmath/amssymb/amsthm/graphicx base
% packages are supplied by the assembler, so the class/externaldocument part of the native preamble is omitted.
% For commutative diagrams we use Xy-pic
\usepackage[all]{xy}

% We use 2cell for 2-commutative diagrams.
\xyoption{2cell}
% Theorem environments.
%
\theoremstyle{plain}
\newtheorem{theorem}[subsection]{Theorem}
\newtheorem{proposition}[subsection]{Proposition}
\newtheorem{lemma}[subsection]{Lemma}

\theoremstyle{definition}
\newtheorem{definition}[subsection]{Definition}
\newtheorem{example}[subsection]{Example}
\newtheorem{exercise}[subsection]{Exercise}
\newtheorem{situation}[subsection]{Situation}

\theoremstyle{remark}
\newtheorem{remark}[subsection]{Remark}
\newtheorem{remarks}[subsection]{Remarks}

\numberwithin{equation}{subsection}

% Macros
%
\def\lim{\mathop{\mathrm{lim}}\nolimits}
\def\colim{\mathop{\mathrm{colim}}\nolimits}
\def\Spec{\mathop{\mathrm{Spec}}}
\def\Hom{\mathop{\mathrm{Hom}}\nolimits}
\def\Ext{\mathop{\mathrm{Ext}}\nolimits}
\def\SheafHom{\mathop{\mathcal{H}\!\mathit{om}}\nolimits}
\def\SheafExt{\mathop{\mathcal{E}\!\mathit{xt}}\nolimits}
\def\Sch{\mathit{Sch}}
\def\Mor{\mathop{\mathrm{Mor}}\nolimits}
\def\Ob{\mathop{\mathrm{Ob}}\nolimits}
\def\Sh{\mathop{\mathit{Sh}}\nolimits}
\def\NL{\mathop{N\!L}\nolimits}
\def\CH{\mathop{\mathrm{CH}}\nolimits}
\def\proetale{{pro\text{-}\acute{e}tale}}
\def\etale{{\acute{e}tale}}
\def\QCoh{\mathit{QCoh}}
\def\Ker{\mathop{\mathrm{Ker}}}
\def\Im{\mathop{\mathrm{Im}}}
\def\Coker{\mathop{\mathrm{Coker}}}
\def\Coim{\mathop{\mathrm{Coim}}}

% Boxtimes
%
\DeclareMathSymbol{\boxtimes}{\mathbin}{AMSa}{"02}

%
% Macros for moduli stacks/spaces
%
\def\QCohstack{\mathcal{QC}\!\mathit{oh}}
\def\Cohstack{\mathcal{C}\!\mathit{oh}}
\def\Spacesstack{\mathcal{S}\!\mathit{paces}}
\def\Quotfunctor{\mathrm{Quot}}
\def\Hilbfunctor{\mathrm{Hilb}}
\def\Curvesstack{\mathcal{C}\!\mathit{urves}}
\def\Polarizedstack{\mathcal{P}\!\mathit{olarized}}
\def\Complexesstack{\mathcal{C}\!\mathit{omplexes}}
% \Pic is the operator that assigns to X its picard group, usage \Pic(X)
% \Picardstack_{X/B} denotes the Picard stack of X over B
% \Picardfunctor_{X/B} denotes the Picard functor of X over B
\def\Pic{\mathop{\mathrm{Pic}}\nolimits}
\def\Picardstack{\mathcal{P}\!\mathit{ic}}
\def\Picardfunctor{\mathrm{Pic}}
\def\Deformationcategory{\mathcal{D}\!\mathit{ef}}
\usepackage{cleveref}
\crefname{originallemma}{original lemma}{original lemmas}
\Crefname{originallemma}{Original lemma}{Original lemmas}
\makeatletter
\newlabel{lemma-finite-flat-weak-assassin-up-down}{{05FN}{source}{}{sourcepointer-1}{}}
\newlabel{lemma-finite-flat-weak-assassin-up-down@cref}{{[originallemma][0][]05FN}{[1][0][]source}{}{}{}}
\makeatother
\crefname{originallemma}{original lemma}{original lemmas}
\Crefname{originallemma}{Original lemma}{Original lemmas}
\makeatletter
\newlabel{morphisms-lemma-etale-locally-standard-etale}{{02GT}{source}{}{sourcepointer-2}{}}
\newlabel{morphisms-lemma-etale-locally-standard-etale@cref}{{[originallemma][0][]02GT}{[1][0][]source}{}{}{}}
\makeatother
\title{Weakly associated points of a quasi-coherent module are detected after an etale morphism}
\author{The Stacks Project Authors}
\date{stacks-project:flat.tex:local-snapshot-20261011; GNU FDL 1.2 or later}
\begin{document}
\maketitle
\noindent\textit{Editorial note.} All delivered mathematical text is the author's own, copied from the pinned Stacks Project chapter file flat.tex; no statement or proof step was rewritten, reordered or removed.\par\medskip
\noindent\textit{Editorial note.} The preamble is the chapter's own retained minimal preamble (sources/stacks-flat/retained-preamble.tex), a verbatim excerpt of the native Xy-pic setup and of the native theorem declarations and mathematical macros; no macro was invented or redefined.\par\medskip
\noindent\textit{Editorial note.} The same-chapter referenced result is located at its \textbackslash{}label line inside the pinned chapter file; the cross-chapter standard-etale lemma is located in the project's own tag list staged in the snapshot (sources/stacks-chow/tags.txt), at the verbatim line carrying its canonical tag.\par\medskip
\noindent\textit{Allowed standard prerequisites (editorial).}\begin{itemize}
\item Weakly associated points of a quasi-coherent module. Introduced earlier in the same project (Divisors, Definition divisors-definition-weakly-associated) and imported by this chapter. Supplies the notion WeakAss appearing on both sides of the asserted equivalence; the reduction that transports it through the etale morphism stays inside the counted proof span.
\item Reduction of an etale morphism to a standard etale morphism. Morphisms, Lemma labelled morphisms-lemma-etale-locally-standard-etale (canonical Stacks tag 02GT); invoked at proof line 367. Supplies the local normal form S = Spec(A), U = Spec(A[x, 1/g]/(f)) with f monic and f' invertible; the identification of that algebra as a localisation of a finite locally free A-algebra is the counted step.
\end{itemize}\medskip
\ifdefined\hypertarget\hypertarget{sourcepointer-1}{}\fi
\noindent\textit{Source reference.} Original lemma 05FN: sources/stacks-flat/pinned-original-source.tex, source lines 312--312. Same-chapter pointer invoked at proof line 375 for the finite locally free case; the two counted steps that put U inside a finite locally free S-scheme (localisation of A[x]/(f) and the standard etale normal form) are retained verbatim in the delivered body.\par\medskip
\ifdefined\hypertarget\hypertarget{sourcepointer-2}{}\fi
\noindent\textit{Source reference.} Original lemma 02GT: sources/stacks-chow/tags.txt, source lines 3240--3240. Cross-chapter standard prerequisite invoked once at proof line 367 for the local normal form of an etale morphism; it supplies the shape S = Spec(A), X = Spec(A[x, 1/g]/(f)) only, and the finite-free-envelope argument is retained in the counted span.\par\medskip
\begin{lemma}
\label{lemma-etale-weak-assassin-up-down}
Let $h : U \to S$ be an \'etale morphism of schemes.
Let $\mathcal{G}$ be a quasi-coherent $\mathcal{O}_S$-module.
Let $u \in U$ be a point with image $s \in S$. Then
$$
u \in \text{WeakAss}(h^*\mathcal{G})
\Leftrightarrow
s \in \text{WeakAss}(\mathcal{G})
$$
\end{lemma}

\begin{proof}
After replacing $S$ and $U$ by affine neighbourhoods of $s$ and $u$
we may assume that $g$ is a standard \'etale morphism of affines, see
Morphisms, Lemma \ref{morphisms-lemma-etale-locally-standard-etale}.
Thus we may assume $S = \Spec(A)$ and
$X = \Spec(A[x, 1/g]/(f))$, where $f$ is monic and $f'$
is invertible in $A[x, 1/g]$.
Note that $A[x, 1/g]/(f) = (A[x]/(f))_g$ is also the localization
of the finite free $A$-algebra $A[x]/(f)$. Hence we may think of
$U$ as an open subscheme of the scheme $T = \Spec(A[x]/(f))$
which is finite locally free over $S$. This reduces us to
Lemma \ref{lemma-finite-flat-weak-assassin-up-down}
above.
\end{proof}

\end{document}
